\subsection{\texorpdfstring{\(\log_a x\) 的导数}{以 a 为底的 x 的对数之导数}}

由于 \(a^x\) 在 \(\mathbf{R}\) 上严格单调（当 \(a \neq 1\) 时），应用反函数求导法则（I, p.~17, prop. 6）便得，对一切 \(x > 0\)

\[
\mathrm{D} \big (\log_ {a} x \big) = \frac {1}{x \log a}\tag{9}
\]

特别地，

\[
\mathrm{D} (\log x) = \frac {1}{x}.\tag{10}
\]

若 u 是在点 \(x_{0}\) 处有导数且满足 \(u(x_{0}) > 0\) 的实函数，则函数 \(\log u\) 有导数，等于 \(u'(x_{0})/u(x_{0})\) （在点 \(x_{0}\) 处）。特别地，有 \(\mathrm{D}\left(\log |x|\right) = 1/|x| = 1/x\) （当 x \textgreater{} 0 时），而

\[
\mathrm{D} \big (\log | x | \big) = - \frac {1}{| x |} = \frac {1}{x}
\]

当 x \textless{} 0 时有上式；换句话说，\(\mathrm{D}\big(\log|x|\big) = 1/x\) 对任何 \(x \neq 0\) 成立。由此得出：若在区间 I 上实函数 u 不为零且有有限导数，则 \(\log|u(x)|\) 在 I 上有导数，等于 \(u'/u\) ；这个导数称为 u 的对数导数。显然，\(|u|^{\alpha}\) 的对数导数是 \(\alpha u'/u\) ，而乘积的对数导数等于各因子的对数导数之和；这些法则常常是计算函数导数最快的方法。特别地，它们再次给出公式

\[
\mathrm{D} \left(x ^ {\alpha}\right) = \alpha x ^ {\alpha - 1}
\]

\[
(\alpha \text {   an   arbitrary   real   number,   } x > 0)\tag{11}
\]

它已用另一种方法证明过（II, p.~69）。

例。若 u 是区间 I 上 \(\neq 0\) 的函数，而 v 是任意实函数，则 \(\log(|u|^{v}) = v \cdot \log |u|\) ，于是当 u 与 v 都可导时

\[
\frac {1}{| u | ^ {v}} \mathrm{D} (| u | ^ {v}) = v ^ {\prime} \log | u | + v \frac {u ^ {\prime}}{u}.
\]

